% chapters/09/contoh.tex
\section{Dua Contoh Terapan}
\label{sec:contoh-bab-9}

\subsection{Contoh Pertama: Melacak Keputusan pada Beberapa Bingkai}

\begin{contohbox}
Sebuah NPC memiliki jari-jari persepsi 260 piksel, jangkauan serang 46 piksel,
pewaktu berdiam 1,5 detik, dan pewaktu tenggang 1 detik. Tabel~\ref{tab:jejak-keputusan}
memperlihatkan keputusannya pada beberapa bingkai yang penting, beserta alasan
masing-masing.
\end{contohbox}

\begin{table}[htbp]
\centering
\caption{Jejak keputusan NPC pada bingkai-bingkai yang penting}
\label{tab:jejak-keputusan}
{\setlength{\tabcolsep}{4pt}\footnotesize%
\begin{tabularx}{\textwidth}{@{}C{1.7cm} C{1.8cm} C{1.7cm} C{1.6cm} C{2.0cm} Y@{}}
\toprule
\headrow \textbf{Bingkai} & \textbf{Keadaan} & \textbf{Pewaktu} & \textbf{Jarak} & \textbf{Keputusan} & \textbf{Alasan} \\
\midrule
90 & berdiam & diam 1,50 & 463 & berpatroli &
Pewaktu berdiam telah mencapai 1,5 detik, sedangkan pemain tidak terlihat. \\
\addlinespace
150 & berpatroli & diam 0,00 & 207 & mengejar &
Pemain masuk ke dalam jari-jari persepsi. \\
\addlinespace
216 & mengejar & hilang 0,00 & 45 & menyerang &
Jarak telah mencapai jangkauan serang 46 piksel. \\
\addlinespace
300 & menyerang & serang 0,10 & 35 & menyerang &
Jarak masih di dalam jangkauan serang, sehingga keadaan tidak berubah. \\
\addlinespace
420 & menyerang & serang 0,70 & 508 & mengejar &
Pemain menjauh melampaui jangkauan serang dikalikan 1,2. \\
\addlinespace
481 & mengejar & hilang 1,00 & 355 & berdiam &
Pewaktu tenggang habis, yaitu satu detik tanpa melihat pemain. \\
\bottomrule
\end{tabularx}}
\end{table}

Perhatikan baris terakhir. Jarak pada bingkai itu 355 piksel, yaitu masih di
luar jari-jari persepsi 260 piksel, sehingga pewaktu tenggang berjalan sampai
habis dan NPC menyerah. Apabila jaraknya kurang dari 260 piksel, NPC akan
melihat pemain kembali, pewaktu tenggangnya dikembalikan ke nol, dan ia terus
mengejar.

Pemeriksaan jejak semacam ini merupakan cara yang paling cepat untuk menemukan
kesalahan pada mesin keadaan. Apabila sebuah perpindahan terjadi terlalu cepat
atau terlambat, penyebabnya akan tampak pada kolom pewaktu.

\subsection{Contoh Kedua: Menghitung Jumlah Serangan}

\begin{contohbox}
Pada skenario pengujian, NPC berada pada keadaan menyerang sejak bingkai 217
sampai bingkai 421, yaitu selama 204 bingkai. Pada laju 60 bingkai per detik,
lamanya adalah

\[
204 \div 60 = 3{,}4 \text{ detik}.
\]

Pewaktu jeda serangan bernilai 0,8 detik, dan serangan pertama terjadi segera
setelah keadaan dimasuki. Banyaknya serangan karena itu adalah

\[
\frac{3{,}4}{0{,}8} + 1 \approx 4{,}25 + 1 \approx 5
\]

apabila dibulatkan ke bawah, yaitu 5 serangan. Program mencatat
\textbf{5 serangan}, sesuai dengan hitungan itu.
\end{contohbox}

Kesesuaian antara hitungan dan catatan program itu berguna untuk dua hal.
Pertama, ia menunjukkan bahwa pewaktu jeda bekerja sebagaimana mestinya.
Kedua, ia memberi Anda cara untuk merancang: apabila Anda menghendaki jumlah
serangan tertentu selama pemain berada di dalam jangkauan, nilai pewaktu jeda
dapat dihitung terbalik dari keinginan itu.

\begin{catatabox}
Hitungan semacam ini hanya berlaku apabila keadaan menyerang
berlangsung tanpa terputus. Apabila pemain keluar masuk jangkauan serang
beberapa kali, jumlah serangan menjadi lebih banyak daripada hitungan itu,
sebab serangan pertama terjadi setiap kali keadaan dimasuki. Kebiasaan
memeriksa ulang hitungan semacam ini terhadap catatan program merupakan bagian
dari kebiasaan mengukur yang telah dibangun sejak Bab~IV.
\end{catatabox}
